% !Mode:: "TeX:UTF-8"

\hitsetup{
  info={
    %******************************
    % 注意：
    %   1. 配置里面不要出现空行
    %   2. 不需要的配置信息可以删除
    %   3. 中英文两版的字段收在 zh={} 与 en={} 里，字段名不带后缀。
    %      这是推荐写法；带后缀的 title-zh={...} 也认，两种等价
    %******************************
    %
    %===========
    % 本科论⽂类型
    %===========
    % 毕业论文还是毕业设计写在 \documentclass 的选项里：
    % degree-level=bachelor 勾毕业论文，degree-level=bachelor* 勾毕业设计。
    %
    %=====
    % 秘级
    %=====
    secret-level={公开},
    classified-index-national={TM301.2},
    classified-index-international={62-5},
    %
    %===========
    % 不分语言的
    %===========
    report-number={no9527}, %编号
    %
    %=========
    % 中文信息
    %=========
    zh={
      % 花括号分组给「题目栏是两条固定线」的封面分行（本科第二封面、博后第一封面）；
      % \\ 给自由排的封面分行；方括号是副标题，写了就排。
      % 带副标题写成 title={...}{...}[一条副标题]。副标题只排在自由排的封面上，
      % 题目栏是两条固定线的那几种（博后第一、威海本科第二、深圳博士中期）没有这一栏
      discipline={工学},
      associate-supervisor={某某某教授}, % 副指导老师
      co-supervisor={某某某教授}, % 联合指导老师
      % 首页最下方那行日期不写就跟着上面的 date 走。要单独指定就取消下一行，
      % 写法与 date 相同。
      % cover-date={2022-06},
      % 日期不写就取编译当天。要指定就写 ISO 数字：yyyy、yyyy-mm、yyyy-mm-dd
      % 都认，月日补不补零都行；中文封面出「2026 年 6 月」，英文封面出
      % 「June, 2026」。写成别的样子就原样排出来，日志里会说一句。
      % date={2026-06},
      % 学生类别不用写，由文档类选项 degree-type 给：academic 出「学术学位论文」，
      % professional 出「专业学位论文」。要改这两条文本本身，用 const 族的
      % student-type-academic 与 student-type-professional。
      % 只有两条都不合适时才写下面这个字段，写了就以写的为准。外面那对括号由
      % 模板加，值里不要自己带。非全日制教育申请学位者写同等学力人员、工程硕士、
      % 工商管理硕士、高级管理人员工商管理硕士、公共管理硕士、中职教师、高校教师等
      % student-type={工程硕士},
      station={哈铁西站}, %博士后站名称
      % 博后封面的四个日期，同样写 ISO 数字。写到哪一档就排到哪一档：
      % 2050-09 出「2050 年 9 月」，2050-09-10 出「2050 年 9 月 10 日」。
      work-finish-date={2050-09}, % 工作完成日期
      report-submit-date={2050-10}, % 报告提交日期
      research-start-date={3050-09-10}, % 研究工作起始时间
      research-end-date={3090-10-10}, % 研究工作期满时间
      % 关键词用“英文逗号”分割
      keywords={\TeX, \LaTeX, CJK, 嗨！, thesis},
    },
    %
    %=========
    % 英文信息
    %=========
    en={
      % 带副标题写成 title={...}[This is the sub title]
      title={Research on key technologies of partial porous externally pressurized gas bearing},
      discipline={Engineering},
      major={Computer Science and Technology},
      affiliation={\emultiline[t]{School of Mechatronics Engineering \\ Mechatronics Engineering}},
      author={Yu Dongmei},
      supervisor={Professor XXX},
      associate-supervisor={XXX},
      % 不写就取编译当天。要指定就写 ISO 数字，模板按各处的语言和精度排：
      % 中文封面出「2017 年 12 月」，英文封面出「December, 2017」。
      % yyyy、yyyy-mm、yyyy-mm-dd 都认，月日补不补零都行。
      date={2017-12},
      % 英文那行同理，不写就由 degree-type 给「Academic Degree」或
      %「Professional Degree」。要写具体的就取消注释。
      % student-type={Master of Art},
      keywords={\TeX, \LaTeX, CJK, template, thesis},
    },
  },
}
